%@ pages 108-124
%@ status translated
% Original section 11, "Catégories de modèles internes", pp. 108--124.
%
% Translation notes (the mathematics is left exactly as printed):
%  * The authors' "cmf" (catégorie de modèles fermée) is rendered "cmc"
%    (closed model category), following the terminology table.
%  * "n-précat (de Segal)" is rendered "Segal n-precategory", as in \S2.
%  * "libre addition de cellules" is rendered "free addition of cells", as
%    in \S5.
%  * "k-groupique" is rendered "k-groupic", as in \S2.
%  * The authors abbreviate "compatible aux changements de base" as "ccb";
%    the translation writes "compatible with base change (cbc)" and uses
%    "cbc" throughout.  The labels \label{ccb1} and \label{cestccb} are the
%    original's and are left unchanged.
%  * In the last paragraph the French sentence breaks off ("...satisfaisant
%    à la propriété universelle homotopique voulue, est obstruée"); the final
%    clause has no subject and is read as saying that the construction is
%    obstructed.  The mathematics is untouched.

\clearpage
\section{Internal model categories}
\label{cmipage}


We want to construct the Segal $(n+1)$-prestack of modules of Segal
$n$-stacks. For this we use the following notion of an {\em internal closed
model category}. Let $M$ be a closed model category. We shall say that $M$ is
{\em internal} if it satisfies the following axioms:
\newline
IM(a)---$M$ admits an internal $Hom$, i.e. the functor $X\mapsto M_1(X\times A,
B)$ is representable by an object $\underline{Hom} (A,B)$;
\newline
IM(b)---the direct product commutes with finite colimits, in particular
$A\times \emptyset = \emptyset$ for the initial object $\emptyset$;
\newline
IM(c)---if $A\rightarrow A'$ and $B\rightarrow B'$ are cofibrations then the
natural morphism
$$
{\rm colim} \left(
\begin{array}{ccc}
A\times B & \rightarrow & A'\times B \\
\downarrow &&  \\
A\times B' & &
\end{array}
\right) \rightarrow A'\times B'
$$
is also a cofibration; and
\newline
IM(d)---weak equivalence is stable under direct product.

These axioms are very close to the axioms of a {\em monoidal closed model
category} in Hovey's sense \cite{HoveyMonMod}. Essentially, one asks that the
direct product define a monoidal structure in the sense of
\cite{HoveyMonMod}, and one asks in addition that the internal $\underline{Hom}$
exist. This last condition is more or less equivalent to the existence of a
certain colimit, but this colimit is perhaps not ``small'' (i.e. indexed by a
set and not a class) and we do not know whether the simple condition of
existence of (small) colimits suffices to guarantee IM(a).

The following proposition recovers the axioms IM1--IM4 of \cite{nCAT} \S 11,
and we regard the axioms IM(a)--IM(d) from now on as constituting the
``official'' version of the notion of internal cmc.

\begin{proposition}
\label{interne}
Let $M$ be an internal closed model category.
\newline
(1) \,\, If $A$ is cofibrant and $B$ fibrant then $\underline{Hom}(A,B)$ is
fibrant; \newline
(2)\,\, If $A\rightarrow A'$ is a cofibration (resp. trivial cofibration) of
cofibrant objects, and if $B'\rightarrow B$ is a fibration (resp. trivial
fibration) between fibrant objects, then the morphism
$$
\underline{Hom}(A',B')\rightarrow
\underline{Hom}(A',B) \times _{\underline{Hom}(A,B) }\underline{Hom}(A,B')
$$
is a fibration (resp. trivial fibration);
\newline
(3)\,\, The internal $\underline{Hom}$ turns every cofibrant coproduct in its
first argument (resp. fibrant fibre product in its second argument) into a
fibrant fibre product.
\end{proposition}
The proof is easy and left to the reader (the arguments are to be found in \S 7
of \cite{nCAT}).
\eop

There is a ``homotopy'' version of the notion of internal $Hom$. Let $A$ be a
simplicial category, and let $x,y\in A$. Suppose that $A$ admits homotopy
direct products; let $h\in A$ and $f: h\times x\rightarrow y$ be a morphism
(i.e. one chooses an object $h\times x$ equipped with morphisms to $h$ and $x$
representing the direct product). One says that $(h,f)$ is a homotopy internal
$Hom$ of $x$ to $y$, if the morphism (essentially well defined)
$$
A_{1/}(u,h)\rightarrow A_{1/}(u\times x, y)
$$
is an equivalence (here again one chooses an object $u\times x$).

\begin{proposition}
\label{hominterne}
Let $M$ be an internal closed model category. Then the simplicial category
$L(M)$ admits an internal $Hom$ in the homotopical sense above.
\end{proposition}
{\em Proof:}
One uses the technique of cofibrant cosimplicial resolutions to compute the
function complexes
(cf \cite{Hirschhorn} \cite{DHK} \cite{DwyerKan3}
\cite{JardineGoerssBook} \cite{Reedy}).
Let us fix $x,y\in M$ with $x$ cofibrant and $y$ fibrant, and denote by
$h:=\underline{Hom}(x,y)$ the internal $\underline{Hom}$ of $M$ (IM(a)). There
is a morphism $h\times x\rightarrow y$. Here and below it is worth remarking
that the direct product of objects of $M$ is also the homotopy direct product
in $L(M)$, by IM(d). Let us fix $u\in M$ and choose a cofibrant Reedy
cosimplicial resolution ${\bf u}\rightarrow u$
(cf \cite{Hirschhorn} \cite{DHK} \cite{DwyerKan3} \cite{JardineGoerssBook}
\cite{Reedy}). One then has the equality between simplicial sets
$$
M_{1/} ({\bf u}, h)= M_{1/}(x\times {\bf u}, y).
$$
Now, since $h$ is fibrant (Proposition \ref{interne} (1)), one also has
(by {\em loc cit.})
$$
M_{1/} ({\bf u}, h)\cong L(M)_{1/}(u,h).
$$
On the other hand, axioms IM(b) and IM(c) imply that $x\times {\bf u}$ is a
cofibrant Reedy cosimplicial resolution of $x\times u$, hence (again by
{\em loc cit.}) one has the equivalence
$$
M_{1/} (x\times {\bf u}, y)\cong L(M)_{1/}(x\times u, y).
$$
These two equivalences imply by definition that $h$ is a homotopy internal
$Hom$ of $x$ to $y$.
\eop

\subnumero{The associated \texorpdfstring{$M_f$}{M\_f}-enriched category}

Let $M$ be an internal closed model category. Let us denote by $\ast$ its
terminal object. One obtains a functor $\sigma : M\rightarrow Sets$ by setting
$\sigma (X):= M_1(\ast , X)$. This construction commutes with direct products.
One has the tautological formula
$$
M_1(A,B) = \sigma \underline{Hom} (A,B).
$$
The internal $\underline{Hom}$ makes it possible to equip $M_{cf}$ with a
structure of category enriched in $M_f$, via the functor $\sigma$. In other
words, for all $A,B\in M_{cf}$ one has $\underline{Hom} (A,B)\in M_f$ and
$M_{cf}(A,B)=\sigma (\underline{Hom} (A,B))$;
for all $A,B,C$ one has a composition
$$
\underline{Hom} (A,B)\times \underline{Hom} (B,C)\rightarrow
\underline{Hom}(A,C)
$$
which is associative (we neglect here the problems connected with the
associativity constraints, which are natural isomorphisms). This composition
is transformed by $\sigma$ into the composition of $M_{cf}$.

If one knows in addition that the $\underline{Hom} (A,B)$ are cofibrant (this
is for example the case when the cofibrations of $M$ are the injections), then
$M_{cf}$ is self-enriched by its internal $Hom$. In particular, the structure
$INT(M_{cf})$ depends only on the category $M_{cf}$ in this case.


One also obtains what one could call a {\em strict $M_f$-Segal category}, and
which we shall denote by $INT(M_{cf})$, by setting
$$
INT(M_{cf})_{p/}(A_0,\ldots , A_p):=
\underline{Hom}(A_0,A_1)\times \ldots \times \underline{Hom}(A_{p-1}, A_p).
$$
More precisely, this object consists of a set $INT(M_{cf})_0$ of objects, and
for every sequence $A_0,\ldots , A_p$ of objects, an object
$INT(M_{cf})_{p/}(A_0,\ldots , A_p)$ of $M_f$ (with the usual morphisms of
functoriality) such that the Segal morphisms are isomorphisms. This is of
course equivalent to the data of an $M_f$-enriched category, and we shall use
the notation $INT(M_{cf})$ indifferently for the $M_f$-enriched category or
for the strict $M_f$-Segal category.



\subnumero{Internal structure for Segal \texorpdfstring{$n$}{n}-categories}



In the case where $M$ is a category of presheaves of sets on some category, and
if the cofibrations are the injective morphisms (or even only injective over a
subset of objects of the base category) then axioms IM(a), IM(b) and IM(c) are
automatically satisfied.
Such is the case for the closed model category of $n$-precategories of
\cite{nCAT} or its variant for Segal $n$-precategories defined earlier in the
present work \ref{cmf}. Such is also the case for the closed model category
of Segal $n$-prestacks \ref{SeCmf}, since a Segal $n$-prestack is nothing other
than a presheaf of sets on $\Theta ^{n+1} \times \Xx$ (and the cofibrations are
the injections of presheaves).
The only difficulty in these examples consists in verifying IM(d).

This has been done
for $n$-precategories in \cite{nCAT} Theorem 5.1. As for the existence of the
closed model category structure, one can either copy the proof, or use this
result via $A\mapsto \Pi _m \circ A$, to obtain the same result for Segal
$n$-precategories. One obtains:

\begin{lemma}
The cmc $nSePC$ of Segal $n$-precategories of \ref{cmf} is
internal.
\end{lemma}
\eop





If one applies the preceding discussion to the closed model category $nSePC$
of Segal $n$-precategories one obtains a strict $1$-category
$INT(nSePC_{cf})$ enriched over $nSePC_f$; and by taking the nerve of this one
obtains an $nSePC_f$-Segal category, i.e. a Segal $n+1$-category which we
denote
$$
nSeCAT:= INT(nSePC_{cf}).
$$
Let us insist on the fact that the objects of $nSeCAT$ are the fibrant Segal
$n$-categories, and for every sequence of objects $A_0,\ldots , A_p$ one has
$$
nSeCAT_{p/}(A_0,\ldots , A_p):= \underline{Hom}(A_0, A_1)\times \ldots
\times \underline{Hom}(A_{p-1}, A_p).
$$
The same construction for the $n+1$-category $nCAT$ of fibrant $n$-categories
was given in \cite{nCAT} \S 7.

We observe that $nSeCAT$ (resp. $nCAT$) is not fibrant, and one will often use
a fibrant replacement which we shall denote $nSeCAT'$ (resp. $nCAT'$).

{\em Set-theoretic remark:}
The objects of $(nSeCAT)_0$ form a class, and technically speaking, $nSeCAT$ is
not an $n+1$-category. To get around this problem, we agree to fix a set (a
very large one!) and to retain as objects of $nSeCAT$ only those $n$-categories
whose underlying set (i.e. union of the $A_M$) is contained in this set. With
this convention, the fibrant replacement $nSeCAT'$ exists. The reader will
verify that our constructions do not take us outside this framework: when one
speaks of limits or colimits, for example, they are always taken along a set of
indices whose cardinality one can bound in advance; it therefore suffices to
choose our initial set large enough for all the limits and colimits that we are
going to take to exist; similarly, when one uses the small object argument, one
always knows in advance a transfinite bound on the number of steps required.




\subnumero{Internal structure for Segal \texorpdfstring{$n$}{n}-stacks}

We shall now adapt to Segal $n$-stacks the construction seen for Segal
$n$-categories.


\begin{lemma}
The cmc $nSePCh$ of Segal $n$-prestacks on a site $\Xx$, is internal.
\end{lemma}
{\em Proof:} as remarked above, the only
difficulty concerns IM(d),
which one proves by induction on $n$. For $n=0$ the property in question
follows easily from the definition of Illusie's equivalence (noting that
passage to the associated sheaf is compatible with the direct product of
presheaves of sets, see for example \cite{Rezk}). We therefore suppose
$n\geq 1$ and the result known for $n-1$. Let $f:A\rightarrow A'$ be a
$\Gg$-equivalence of Segal $n$-prestacks, and let $B$ be a Segal $n$-prestack.
One may suppose (thanks to the corresponding result for Segal $n$-precategories)
that $A(X)$, $A'(X)$ and $B(X)$ are Segal $n$-categories. One has
$$
\tau _{\leq 0} (A\times B)= \tau _{\leq 0}(A)\times \tau _{\leq 0} (B)
$$
(and likewise for $A'$)
hence the morphism
$$
\tau _{\leq 0} (A\times B)\rightarrow \tau _{\leq 0} (A'\times B)
$$
induces a surjection on the associated sheaves of sets. On the other hand,
for $(x,a)$ and $(y,b)$ in $(A\times B)_0(X)$, one has
$$
(A\times B)_{1/}((x,a), (y,b))= A_{1/}(x,y)\times B_{1/}(a,b)
$$
and likewise for $A'$; one therefore obtains (from the result for $n-1$) that
the morphism induced by $f$
$$
(A\times B)_{1/}((x,a), (y,b))\rightarrow
(A'\times B)_{1/}((fx,a), (fy,b))
$$
is a $\Gg$-equivalence. By definition (see \S 3) this means that $f$ is a
$\Gg$-equivalence. This gives IM(d).
\eop



Now,
starting from $M=nSePCh$, one can construct the Segal $n+1$-prestack of fibrant
Segal $n$-stacks, denoted
$$
nSe\underline{CHAMP}(\Xx ):= INT(nSePCh_{cf}).
$$
The objects of $nSe\underline{CHAMP}(\Xx )$ over $X\in \Xx$ are the fibrant
Segal $n$-prestacks on $\Xx /X$. If $A_0,\ldots , A_p$ are such objects, then
one defines on $\Xx /X$ the Segal $n$-prestack
$$
nSe\underline{CHAMP}(\Xx )_{p/}(A_0,\ldots , A_p)=
\underline{Hom}(A_0,A_1)\times \ldots \times \underline{Hom}(A_{p-1}, A_p).
$$

By applying the global sections functor $\Gamma$ to
$nSe\underline{CHAMP}(\Xx )$, one obtains {\em the Segal $n+1$-category of
Segal $n$-stacks on $\Xx$}, which one could denote
$$
nSeCHAMP(\Xx ).
$$
In fact, if one analyses the preceding sentence, one finds that
the objects of $nSeCHAMP(\Xx )$
are the Segal $n$-prestacks $A$ such that $A|_{\Xx /X}$ is
fibrant for all $X\in \Xx$. If $\Xx$ has no terminal object, this condition
may be different from the (stronger) condition that $A$ be fibrant
(note the counterexample in \S 6),
and we shall modify the definition slightly so as to retain in
$nSeCHAMP(\Xx )$ only those objects which are fibrant over $\Xx$.

In conclusion, one prefers to denote by
$nSeCHAMP(\Xx )$
the $n+1$-category whose objects
are the $\Gg$-fibrant Segal $n$-prestacks on the site
$\Xx$. If $A_0,\ldots , A_p$ are
such objects, then one takes as previously for
$nSeCHAMP(\Xx )_{p/}(A_0,\ldots , A_p)$
the product of $n$-precategories
$\Gamma \underline{Hom}(A_0,A_1)\times \ldots \times
\Gamma \underline{Hom}(A_{p-1}, A_p).$

It is worth recalling, in order to understand this definition, that
$\Gamma \underline{Hom}(A,B)$ can be defined as the representative of the
functor
$nSePC\rightarrow Ens$
$$
E\mapsto Hom _{nSePCh}(A\times \underline{E},B),
$$
where $\underline{E}$ is the constant Segal $n$-prestack with values $E$.


It is immediate that for every $X\in \Xx$, one has
$$
nSe\underline{CHAMP}(\Xx )(X) = nSeCHAMP(\Xx /X),
$$
and this could serve as a definition of $nSe\underline{CHAMP}(\Xx )$.

Finally, one can make here the same ``set-theoretic remark''
as above (which explains
why one allows oneself to take fibrant replacements
$nSe\underline{CHAMP}(\Xx )'$).

\subnumero{Computation with the HBKQ-type structure}

We indicate here how to use the structure of theorem
\ref{cmfHirschho} to compute $\Gamma \underline{Hom}$ and
$\underline{Hom}$ (the latter in the case where $\Xx$ admits fibre
products).

\begin{lemma}
\label{hbkqinterne}
If the site $\Xx$ admits direct products, then the cmc of HBKQ type of
\ref{cmfHirschho} is internal.
\end{lemma}
{\em Proof:} It suffices to remark
that the product of two free additions of cells over $X$
and over $Y$, is a free addition over $X\times Y$.
\eop

{\em Remark:} If one wishes to apply this lemma on the sites $\Xx /X$ then
these must admit direct products, i.e. $\Xx$ must admit fibre products.

In view of this lemma, one could, in the case where $\Xx$ admits direct
products, define an HBKQ version of $nSeCHAMP(\Xx )$ and
$nSe\underline{CHAMP}(\Xx )$. We point out now that one would obtain
essentially the same objects as above.

We shall begin by showing a property valid even if $\Xx$ does not admit
all direct products.


Let $E$ be a Segal $n$-precategory, and denote by
$\underline{E}$ the constant Segal $n$-prestack on $\Xx$ with value $E$.
If $A\rightarrow A'$ is the cofibration freely generated by the addition
of a cell over $X\in \Xx$ (i.e. generated by a cofibration of
Segal $n$-precategories $A(X)\rightarrow C$) then $A\times
\underline{E}\rightarrow A'\times \underline{E}$ is also freely generated
by addition of the cell $A(X)\times E \rightarrow C\times E$.
It follows (from the fact that the direct product with $\underline{E}$ preserves
transfinite compositions and retractions) that if
$$
A\rightarrow A'
$$
is a cofibration for the structure of HBKQ, then
the same is true of
$$
A\times
\underline{E}\rightarrow A'\times \underline{E}.
$$
Let now $A$ be cofibrant and $B$ fibrant for the structure of
HBKQ (Theorem \ref{cmfHirschho}). The Segal $n$-precategory
$\Gamma \underline{Hom}(A,B)$ represents the functor $nSePC\rightarrow Ens$
$$
E\mapsto Hom _{nSePCh}(A\times \underline{E}, B).
$$
The fact that the $A\times \underline{E}$ are cofibrant (for
\ref{cmfHirschho})
implies that if $B\rightarrow B'$ is an equivalence between two
Segal $n$-prestacks fibrant for \ref{cmfHirschho}, then the induced
morphism
$$
\Gamma \underline{Hom}(A,B)\rightarrow \Gamma \underline{Hom}(A,B')
$$
is an equivalence of Segal $n$-precategories (both are fibrant
Segal $n$-precategories). In particular
one may choose $B'$
fibrant with respect to the structure of theorem \ref{cmf},
in which case one obtains the
``right'' Segal $n$-category of morphisms from $A$ to $B$. The conclusion
is the following lemma, which says that $A$ and $B$ may be used to
compute the Segal $n$-categories of morphisms in $nSeCHAMP (\Xx )$.

Let us note, before stating the lemma, that $B$ is fibrant for the structure
\ref{cmfHirschho} if and only if for all $X\in \Xx$, $B(X)$ is fibrant.
This is therefore a pointwise condition (satisfied for example by every
functor $\Xx ^o \rightarrow nSeCAT$).

\begin{lemma}
\label{bonhomHirschho}
Let $A$ be cofibrant and $B$ fibrant for the structure of
HBKQ (Theorem \ref{cmfHirschho}). Let $A'$ and $B'$ be
their fibrant replacements for the structure of \ref{cmf}. Then
there is an equivalence of Segal $n$-categories
$$
nSeCHAMP(\Xx )_{1/}(A',B')=\Gamma \underline{Hom}(A',B')\cong
\Gamma \underline{Hom}(A,B).
$$
\end{lemma}
\eop

We point out that one can compute in the same way the Segal $n$-prestacks
of morphisms in $nSe\underline{CHAMP}(\Xx )$, {\em provided
that $\Xx$ admits fibre products}; thanks to lemma
\ref{hbkqinterne}:

\begin{lemma}
\label{bonhomHirschho2}
Suppose that the site $\Xx$ admits fibre products.
Let $A$ be cofibrant and $B$ fibrant for the structure of
HBKQ (Theorem \ref{cmfHirschho}). Let $A'$ and $B'$ be
their fibrant replacements for the structure of \ref{cmf}. Then
there is an equivalence of Segal $n$-categories
$$
nSe\underline{CHAMP}(\Xx )_{1/}(A',B')=\underline{Hom}(A',B')\cong
\underline{Hom}(A,B).
$$
\end{lemma}
\eop

Lemma \ref{bonhomHirschho} helps us to make explicit the notion of ``descent
data''. Let $X\in \Xx$ and let $\Bb \subset \Xx /X$ be a sieve
covering $X$. Let
$$
{\cal D}_{\Bb} \rightarrow \ast _{\Bb}
$$
be the morphism of Segal $n$-prestacks (which may be regarded as
being over $\Xx /X$ or $\Bb$) constructed in \S 5.
Let $A$ be a Segal $n$-prestack on $\Xx$ fibrant for the structure
\ref{cmfHirschho}, i.e. a presheaf on $\Xx$ with values in the
category
$nSePC_f$ of fibrant Segal $n$-categories.
Let $A\rightarrow A'$ be a fibrant replacement for the structure of \ref{cmf}
for the coarse topology. Then one has the equivalence
$$
\Gamma \underline{Hom} (\ast _{\Bb}, A'|_{\Bb})
\cong \Gamma \underline{Hom}({\cal D}_{\Bb}, A|_{\Bb}).
$$
In particular, taking $\tau _{\leq 0}$,
one obtains
the following corollary.

\begin{corollary}
\label{calculDescente}
Let $A$ be a Segal $n$-prestack on $\Xx$ fibrant for the structure
\ref{cmfHirschho}, i.e. a presheaf on $\Xx$ with values in the
category
$nSePC_f$ of fibrant Segal $n$-categories.
Let $X\in \Xx$ and  $\Bb \subset \Xx /X$ a sieve covering $X$, and let
$$
{\cal D}_{\Bb} \rightarrow \ast _{\Bb}
$$
be the ``resolution'' constructed in \S 5. Then the set
``of descent data''
$$
[\ast _{\Bb}, A|_{\Bb}]_{\Bb ^{\rm
gro}}
$$
is computed as the set of morphisms (in $nSePCh$) from ${\cal
D}_{\Bb}$ to
$A$, modulo the relation which
identifies two morphisms if they are linked by a
morphism ${\cal D}_{\Bb} \times \underline{\overline{I}}\rightarrow A$.
\end{corollary}
{\em Proof:}
Let $A\rightarrow A'$ be a fibrant replacement for the structure of \ref{cmf}
for $\Xx ^{\rm gro}$.
The set of descent data is computed as
$$
[\ast _{\Bb}, A|_{\Bb}]_{\Bb ^{\rm
gro}}
=\tau _{\leq 0}\Gamma \underline{Hom} (\ast _{\Bb}, A'|_{\Bb}).
$$
By the equivalence above this is isomorphic to
$$
\tau _{\leq 0}\Gamma \underline{Hom} ({\cal D}_{\Bb}, A|_{\Bb}).
$$
Since
$\Gamma \underline{Hom} ({\cal D}_{\Bb}, A|_{\Bb})$ is a fibrant Segal
$n$-precategory, its truncation $\tau _{\leq 0}$ is computed as the set of
objects
modulo the relation which identifies
two objects linked by a morphism with source $\overline{I}$.
\eop

\subnumero{Relation with the simplicial structure}

According to Dwyer--Kan \cite{DwyerKan3}, the ``simplicial structure'' of a
closed model category is well defined up to homotopy, and one would like to
compare it with the internal structure when the latter exists.
It would be possible to formulate this question in full generality but
for convenience (and because it suffices for our applications) we shall
only look at the case of an internal closed model category with
compatible simplicial structure. We begin with the
definition: if $M$ is an internal closed model category, a
{\em compatible simplicial structure} is a functor $R: \Delta \rightarrow M$
such that the simplicial category $M_{\ast}$
defined by
$$
Hom _{M_k}(x,y):= Hom _M(R(k) \times x, y)
$$
is a simplicial model category in Quillen's sense.
One can also write
$$
Hom_{M_{\ast}}(x,y) = \left( k\mapsto Hom _M(R(k)\times x, y) \right) .
$$

The functor
$R$ extends, by passage to finite limits, to a functor (which we shall still
denote by $R$) from the category
of simplicial sets having only a finite number of non-degenerate simplices
to $M$. In the notation introduced by Quillen for the simplicial
structure, one has
$$
K\otimes x = R(K)\times x,
$$
and
$$
Hom (K, y) = \underline{Hom}(R(K), y)
$$
($\underline{Hom}$ is the internal $Hom$
of $M$). One can also take these formulas as
an alternative definition of the simplicial structure associated with $R$
(this point of view corresponds to
the definition of simplicial model category
given by Goerss--Jardine \cite{JardineGoerssBook}).

If $(M,R)$ is an internal closed model category with compatible simplicial
structure, one obtains a functor $\Xi : M\rightarrow EnsSpl$
from $M$ to simplicial sets, by setting for $x\in M$,
$$
\Xi (x)_p := Hom _M(R(p), x),
$$
or in terms of representable functors,
$$
Hom _{EnsSpl}(K, \Xi (x)):= Hom _M(R(K), x).
$$
In other words $\Xi (x)$ is the simplicial set
$Hom_{M_{\ast}}(\ast , x)$ where $\ast$ is the terminal object of $M$.
The functor $\Xi$ is compatible with direct products. Consequently,
starting from an $M$-enriched category $C$, one can manufacture the
category $\Xi \circ  C$ which is
enriched over simplicial sets.

\begin{lemma}
\label{splinterne}
Let $(M,R)$ be an internal closed model category with compatible
simplicial structure, and let $INT(M_{cf})$ be the $M_f$-enriched category
constructed above. On the other hand, let $M_{cf,\ast}$ be the simplicial
category of cofibrant and fibrant objects with its simplicial structure
determined by $R$. Then one has an isomorphism of simplicial categories
$$
M_{cf, \ast} = \Xi \circ  INT(M_{cf}).
$$
Consequently, if $L(M)$ is the Dwyer--Kan localisation of $M$ with respect to
the weak equivalences, one has an equivalence of simplicial categories
$$
L(M)\cong \Xi \circ INT(M_{cf}).
$$
\end{lemma}
{\em Proof:}
The objects of $M_{cf, \ast}$ and of $INT(M_{cf})$ are the cofibrant and
fibrant objects of $M$. For two such objects $x,y$, one has by definition
$$
INT(M_{cf})_{1/}(x,y):= \underline{Hom}(x,y),
$$
and therefore
$$
(\Xi \circ INT(M_{cf}))_{1/}(x,y):= \Xi (INT(M_{cf})_{1/}(x,y))
= \left( p\mapsto Hom (R(p), \underline{Hom}(x,y)) \right)
$$
$$
= \left( p\mapsto Hom (R(p)\times x,y) \right)
= Hom _{M_{cf, \ast}}(x,y).
$$
This gives the desired isomorphism. The second part is a consequence
of the equivalence $L(M) \cong M_{cf, \ast}$ see \cite{DwyerKan2}.
\eop

We shall apply all this to Segal $n$-categories and $n$-stacks.
One verifies easily that one can equip the
internal closed model category $nSePC$ of Segal $n$-precategories
with a compatible simplicial structure by taking for $R(p)$
the category $\overline{I}^{(p)}$
having $p+1$ objects $0',\ldots , p'$
and a single isomorphism between each pair of objects. Let
$\Xi : nSePC \rightarrow EnsSpl$ be the corresponding functor.
% The French prints "cat'egorie" here (backslash missing); we print
% "category".

\begin{lemma}
\label{sigma}
If $A$ is a fibrant (and automatically cofibrant) Segal $n$-category then
there is a natural equivalence of simplicial sets
$$
\Xi (A) \stackrel{\cong}{\rightarrow} A^{int, 0}.
$$
\end{lemma}
{\em Proof:}
One has a morphism of Segal $n$-precategories
$$
R\Xi (A) \rightarrow A.
$$
Noting that $SeCat(R\Xi (A))$ is a Segal $n$-groupoid (as a
homotopy colimit of Segal $n$-groupoids), and
using the notation $A^{int, 0'}$ of \S 2, one obtains the equivalence
$$
R\Xi (A)\stackrel{\cong}{\rightarrow}
SeCat(R\Xi (A))^{int, 0'} \rightarrow A^{int, 0'}
$$
(one must choose a retraction $SeCat(A)\rightarrow A$,
which is possible because $A$ is
fibrant). The composite is the morphism above. If $\Re$ is the
realisation as a
simplicial set, one has by definition $A^{int, 0}:= \Re A^{int, 0'}$,
whence a morphism $$
\Re (R\Xi (A)) \rightarrow A^{int, 0}.
$$
If one denotes by $I^{(p)}$ the category with $p+1$ objects $0',\ldots , p'$
and a morphism $i'\rightarrow j'$ for all $i\leq j$, then
the simplicial set
$$
\Re (I^{(p)})= h(p)
$$
is the standard $p$-simplex, and the inclusion
$I^{(p)} \subset \overline{I}^{(p)}$ therefore provides a morphism
$$
h(p)\rightarrow \Re (R(p))
$$
which extends by passage to
colimits to a morphism $K\rightarrow \Re (R(K))$ for every
simplicial set $K$. In particular, one has a morphism
$$
\Xi (A)\rightarrow \Re(R\Xi (A)),
$$
which composed with the morphism above gives
$$
\Xi (A)\rightarrow A^{int, 0}.
$$
We need this to be an equivalence.

The construction $\Xi$ respects the equivalences between fibrant objects.
Using the equality $\Xi (A) = \Xi
(A^{int, 0'})$, one may---after replacing $A$ by $A^{int, 0'}$--
suppose that $A$ is a Segal $n$-groupoid. The
latter corresponds to a
topological space according to the work of
Tamsamani \cite{Tamsamani} (Tamsamani's techniques
apply directly to the ``Segal'' framework because he works in fact
in this framework, and only performs the truncation at the end of the
argument).
One may therefore
suppose that there exists a topological space $Y$ with $A= \Pi _{n,Se}(Y)$. One
notes that the latter is automatically fibrant. One has a weak equivalence of
spaces $|\Re (A)|\rightarrow Y$ (again thanks to \cite{Tamsamani}).
Moreover, the simplicial set
$\Xi (A)$ is a variant of the ``singular complex'' of $Y$, corresponding
to the functor $\Delta \rightarrow Top$ which to $p\in
\Delta$ associates $\Re (R(p))$. It follows that $| \Xi (A)| \rightarrow Y$ is
a weak equivalence.  One obtains that the morphism of topological
spaces
$$
| \Xi (A)|\rightarrow |\Re (A)|
$$
is a weak equivalence, which implies that
$$
\Xi (A)\rightarrow \Re (A) = A^{int, 0}
$$
is an equivalence of simplicial sets.
\eop




\begin{theorem}
\label{intereqloc}
There is a natural equivalence of Segal $1$-categories
$$
nSeCAT ^{int,1} \cong L(nSePC ).
$$
Moreover,
if $\Xx$ is a site, there is a natural equivalence (and functorial in
$\Xx$)
$$
nSeCHAMP (\Xx )^{int, 1} \cong L(nSePCh(\Xx )).
$$
Here the $L(\, )$ are the Dwyer--Kan localisations of the model
categories with respect to the weak equivalences.
\end{theorem}
{\em Proof:}
First $nSeCAT ^{int,1}$ is obtained from $nSeCAT$ by taking
the interior $(\, )^{int, 0}$ of the Segal $n$-categories $nSeCAT _{1/}(A,B)$
which are fibrant.
By lemma \ref{sigma}, one obtains the interior of an object
by applying the construction $\Xi$. One therefore has an equivalence
$$
nSeCAT ^{int,1}\cong \Xi \circ nSeCAT,
$$
and lemma \ref{splinterne} gives the equivalence with
$L(nSePC)$.

The functor $\underline{R}: \Delta \rightarrow nSePCh$ which to $p$
associates the constant presheaf with value $R(p)\in nSePC$,
defines a compatible
simplicial structure on the internal closed model category
$nSePCh$. If one denotes by $\Xi _{\underline{R}}$ the construction $\Xi$
associated with this structure (and by $\Xi _R$ that associated with the
structure $R$ on $nSePC$), one has, for every fibrant Segal $n$-prestack $A$,
$$
\Xi _{\underline{R}}(A) = \Xi _R (\Gamma (\Xx , A)).
$$
Recall that we have set
$$
nSe\underline{CHAMP} := INT (nSePCh_{cf}),
$$
and
$$
nSeCHAMP _{1/} (A, B) := \Gamma (\Xx , nSe\underline{CHAMP}_{1/}(A,B)).
$$
One obtains the equality
$$
(\Xi _{\underline{R}}\circ nSe\underline{CHAMP})_{1/}(A,B)
= (\Xi _R \circ nSeCHAMP )_{1/}(A,B),
$$
in other words
$$
\Xi _{\underline{R}}\circ INT (nSePCh_{cf} )=
\Xi _{\underline{R}}\circ nSe\underline{CHAMP}=
\Xi _R \circ nSeCHAMP.
$$
Lemma \ref{splinterne} then gives the equivalence
$$
L(nSePCh )\cong \Xi _R \circ nSeCHAMP.
$$
As for the first part of the theorem, by
using that the $nSeCHAMP
_{1/}(A,B)$ are fibrant and by applying lemma \ref{sigma} to them, one
obtains the equivalence
$$
\Xi _R \circ nSeCHAMP\cong  nSeCHAMP ^{int, 1}.
$$
These two equivalences give the second part of the theorem.
\eop


\subnumero{The relative internal \texorpdfstring{$\underline{Hom}$}{Hom}}

We present here an often useful extension of the notion of
internal $\underline{Hom}$: it is the ``fibre-by-fibre'' $\underline{Hom}$ for a morphism
$f:A\rightarrow B$, which we shall denote $\underline{Hom}(A/B, C)\rightarrow B$. It
will exist only under certain conditions on the morphism $f$.

Let $M$ be a category of presheaves (i.e. suppose that there
exists another category $\Yy$ such that $M$ is equivalent to the
category of presheaves of sets on $\Yy$). Let
$A,B,C\in M$ with
an arrow $f:A\rightarrow B$. Let us set, for $E\in M$,
$$
H(E):= \{ E\rightarrow B, \, \, E\times _BA\rightarrow C\} .
$$
We claim that this functor is represented by an object
$\underline{Hom}(A/B,
C)$ of $M$. Indeed, $H(-)$ turns colimits into limits,
which suffices
to guarantee its representability.

Suppose moreover that $M$ is an internal cmc. We shall say that a fibration
$f:A\rightarrow B$ is {\em compatible with base change (cbc)}
if, for
every diagram
$$
B'\stackrel{a}{\rightarrow} B'' \rightarrow B
$$
where $a$ is a weak equivalence, one has that
$$
A\times _BB'\rightarrow A\times _BB''
$$
is also a weak equivalence. One notes that if $M$ is proper
(e.g. for $M=0SePC$
or $M=0SePCh$) then every fibration is compatible with base
change---this is the definition of ``proper''.

\begin{lemma}
\label{ccb1}
If
$$
\begin{array}{ccc}
A&\stackrel{f}{\rightarrow}& B\\
\downarrow && \downarrow \\
A'&\stackrel{f'}{\rightarrow}& B'
\end{array}
$$
is a diagram with $f$ and $f'$ fibrations and the vertical arrows
equivalences, then:
\newline
---if $f'$ is cbc then $f$ is cbc;
\newline
---if $B\rightarrow B'$ is a fibration and $f$ is cbc then $f'$
is cbc; and
\newline
---if $B$ and $B'$ are fibrant and $f$ is cbc then $f'$ is cbc.
\end{lemma}
{\em Proof:}
Left to the reader (in the third case one may choose
two fibrations $B\leftarrow B'' \rightarrow B'$ and
apply the second case).
\eop

The following proposition establishes the relation between the condition
``cbc''
and the relative internal $\underline{Hom}$: for the latter to have a
homotopical meaning the morphism $f:A\rightarrow B$ must be cbc.

\begin{proposition}
\label{hominternerelatif}
Let $M$ be an internal cmc which is a category of presheaves.
Suppose
moreover that $cof \subset M$ is the category of injections of
presheaves. Let
$f:A\rightarrow B$ be a fibration compatible with base change. Then the
functor
$$
C\mapsto \underline{Hom}(A/B, C)
$$
turns fibrations into
fibrations, in particular if $C$ is fibrant then the structure morphism
$\underline{Hom}(A/B, C)\rightarrow B$ is a fibration. This construction
is---up to equivalence---invariant under equivalence in the
variables $A$ and $C$ (these invariances are subject to the condition that the
arrow $f$ remain a cbc fibration and that $C$ remain fibrant).
If
$$
\begin{array}{ccc}
A' &\rightarrow & A \\
\downarrow && \downarrow \\
B'&\rightarrow &B
\end{array}
$$
is a cartesian diagram
then one has the formula
$$
\underline{Hom}(A'/B', C) = \underline{Hom}(A/B,C)\times _BB'.
$$
If the vertical arrows are cbc fibrations and the morphism
$B\rightarrow
B'$ is either a trivial fibration or an equivalence between two
fibrant objects, then the induced morphism
$$
\underline{Hom}(A'/B', C) \rightarrow  \underline{Hom}(A/B,C)
$$
is an equivalence.
\end{proposition}
{\em Proof:} Let $g:C'\rightarrow C$ be a fibration; let us show that
the arrow
$\underline{Hom}(A/B,g)$ satisfies the lifting property for a
trivial cofibration $E\rightarrow E'$. This property amounts to the
lifting property for $g$ with respect to the morphism
$$
j:A\times _BE\rightarrow A\times _BE' .
$$
The fact that $cof$ is compatible with fibre products implies that $j$
is a cofibration; and the cbc property for $f$ implies that $j$
is a
weak equivalence; whence the desired lifting property.

One notes in the same way that if $g:C'\rightarrow C$ is a trivial
fibration then $g$ induces a trivial fibration $\underline{Hom}(A/B,g)$.
This
implies the invariance of the construction under
equivalence in the fibrant argument
$C$ (indeed every equivalence fits into a
triangle---commutative up to a retraction---in which the other two
sides
are trivial fibrations).

Let $A\stackrel{\cong}{\hookrightarrow}A'\rightarrow B$ be a diagram
where the
two morphisms $f:A\rightarrow B$ and $f':A'\rightarrow B$ are
cbc fibrations. In this case, for $C$ fibrant, the morphism
$$
\underline{Hom}(A'/B,C)\rightarrow \underline{Hom}(A/B, C)
$$
is a trivial fibration: for a cofibration $E\rightarrow E'$ it
is a matter of seeing that
$$
A\times _B E'
\cup ^{A\times _BE} A'\times _B E \rightarrow A'\times _BE'
$$
is a trivial cofibration, which is the case thanks to the fact that
$$
A\times _BE\rightarrow A'\times _BE \;\;\; and \;\;\;
A\times _BE' \rightarrow A'\times _BE'
$$
are trivial cofibrations (for this one does not need the
cbc property). It follows that the formation of $\underline{Hom}(A/B, C)$ is
invariant
up to
equivalence under equivalence in the variable $A$ (provided that
$f$ remain fibrant and cbc and that $C$ be fibrant).

Finally, the formula in the case of a base change $B\rightarrow B'$
is an immediate consequence of the definition of the relative internal
$\underline{Hom}$
as a representable functor. If the morphism $B'\rightarrow B$
is a trivial fibration,
the fibre product with it still induces a
trivial fibration. If this arrow is an equivalence between fibrant objects
one may reduce (by introducing a $B"$ as at the beginning of the
proof) to the case of a trivial fibration, which gives the statement
of invariance under change of $B$.
\eop

{\em Question:} If $f:A\rightarrow B$ is cbc, is the same true of
$\underline{Hom}(A/B,C)\rightarrow B$?

\begin{lemma}
\label{propunivHominterne}
Let $M$ be a cmc as in the proposition above. Let $f:A\rightarrow B$
be a fibration compatible with base change and let $C$ be a fibrant object.
For every morphism $B'\rightarrow B$, one has
$$
Sect (B', \underline{Hom}(A/B, C)\times _BB') =
\underline{Hom}(A\times _BB',C).
$$
Here $Sect(U,V)$ denotes,
for a given
morphism $V\rightarrow U$,
the fibre of
$\underline{Hom}(U,V)$ over $1_U\in \underline{Hom}(U,U)$.
\end{lemma}
{\em Proof:}
Use the defining property of $\underline{Hom}(A/B, C)$.
A morphism
$$
E\rightarrow Sect (B', \underline{Hom}(A/B, C)\times _BB')
$$
is a morphism $E\times B'\rightarrow \underline{Hom}(A/B, C)$
whose composite to $B$ factors through $B'$. Or again a
morphism
$$
E\times (B'\times _BA)= (E\times B')\times _B A\rightarrow C,
$$
i.e. a morphism $E\rightarrow \underline{Hom}(A\times _BB',C)$.
\eop


If the answer to the preceding question is affirmative,
or else if $B$
and $B'$ are fibrant, or else if every fibration with base $B$
is cbc (cf. the lemma below for example), then the universal
property
of this lemma is also a homotopy universal property because
the morphisms entering into the fibre
products are fibrant,
the fibre products in question are invariant under equivalence,
and the sections are taken along a fibrant morphism.

We now apply this discussion to our examples $M=nSePC$ and
$M=nSePCh$. In fact it will suffice to treat the example of Segal
$n$-prestacks $M=nSePCh$ because one recovers the case of Segal
$n$-precategories ($M=nSePC$)
as the ``pointwise'' case i.e. the case where the underlying site is
$\Xx =\ast$.

One notes first that if $n=0$ (i.e. in the case of simplicial presheaves)
then $M$ is proper and every fibration $f$ is cbc. For $n>0$ one
obtains that
the objects which are essentially Segal $0$-prestacks, i.e. the
$0$-groupic objects or Segal $n$-groupoids, behave
like
simplicial presheaves. One obtains moreover a small improvement,
namely that it
suffices for the base $B$ to be $0$-groupic.

\begin{lemma}
\label{cestccb}
In the cmc $M=nSePCh$ of Segal $n$-prestacks \ref{SeCmf}, a
fibration
$f:A\rightarrow B$ with base $B$ such that each $B(X)$ is a
$0$-groupic Segal $n$-category (i.e. a Segal $n$-groupoid), is compatible with
base change.
\end{lemma}
{\em Proof:}
One may suppose that $B$ is fibrant (because of the invariance under
change
of $A$).  In view of the partial statement of property \ref{proper},
it suffices
to see the compatibility with base changes of the form
$$
B'\rightarrow SeCat(B')\rightarrow B.
$$
For this it suffices to see the compatibility under the same
base
changes over each object $X\in \Xx$. One is therefore reduced to the pointwise
case
$\Xx =\ast$. Recall that in this case the operation $SeCat$ may be viewed
as a transfinite composite of coproducts with standard cofibrations
denoted $\Sigma \rightarrow h(M)$ in \cite{nCAT} (here $h(M)$ is the
Segal $n$-precategory represented by $M\in \Theta^{n+1}$). One may therefore
reduce
to the case of a base change of the form
$$
\Sigma \rightarrow h(M) \rightarrow B.
$$
Let $\overline{h(M)}$ be the ``$0$-groupic completion'' of $h(M)$,
obtained for
example by taking the $\Pi _{n,Se}$ of the topological realisation
of $h(M)$.
The fact that $B$ is $0$-groupic and fibrant implies the existence of a
factorisation of the form
$$
\Sigma \rightarrow h(M) \rightarrow \overline{h(M)} \rightarrow B.
$$
One may therefore reduce to the case $B= \overline{h(M)}$. Now the
morphism
$B\rightarrow \ast$ is a trivial fibration so that, if
$$
\begin{array}{ccc}
A&\rightarrow &A'\\
\downarrow &&\downarrow \\
B&\rightarrow &\ast
\end{array}
$$
is a factorisation of $A\rightarrow \ast$ with $A'$ fibrant,
the cbc
property for $A\rightarrow B$ is equivalent to this property for
$A'\rightarrow \ast$. But the cbc property for $A'\rightarrow \ast$
follows
from the stability of weak equivalences under direct
product, which is none other than IM(d) for $M=nSePC$ (recall that we referred,
for its proof, to the same property for
non-Segal $n$-precategories, which in turn was one of the
main steps of \cite{nCAT}).
\eop

As a consequence of this lemma, for every fibration $f:A\rightarrow B$
of Segal $n$-prestacks with $B$ $0$-groupic, and for
every
fibrant Segal $n$-prestack $C$, the morphism
$$
\underline{Hom}(A/B,C)\rightarrow B
$$
is a fibration possessing
property \ref{propunivHominterne} which is a homotopy universal property.

The counterexample to properness given in \cite{nCAT} is the
morphism
$$
a_{02}: A=I\rightarrow I^{(2)}=B
$$
where $B$ is the category with three objects $0,1,2$ and arrows
$0\rightarrow 1\rightarrow 2$, and $a_{02}$ is the inclusion of the arrow
$0\rightarrow 2$. This morphism is a fibration but is not compatible with
base change. This example shows that a relative internal
$\underline{Hom}$ of the form $\underline{Hom}(A/B,C)$
satisfying the desired homotopy universal property does not always exist;
the construction is obstructed.
